\documentclass[12pt, a4paper]{article}

% Paquetes requeridos para el circuito
\usepackage[utf8]{inputenc}
\usepackage[T1]{fontenc}
\usepackage{circuitikz} 

\begin{document}

\begin{center}
    \begin{circuitikz}[american, scale=1.0, transform shape]
        
        % 1. Raíl inferior común (nodo 'b')
        \draw (0,0) -- (8,0) node[circ] {} node[right] {b};
        
        % 2. Fuente independiente de tensión V1 entre 'b' y 'c'
        \draw (0,0) node[circ] {} to[V, l=$V_1$] (0,3) node[circ] {} node[above] {c};
        
        % 3. Resistencia R1 entre 'c' y 'd'
        \draw (0,3) to[R, l=$R_1$] (4,3);
        
        % 4. Nodo 'd' con la tensión de control vx (tensión de 'd' respecto a 'b')
        \node[circ] at (4,3) (d) {};
        \node[above] at (d) {$d$};
        \draw (d) to[open, v=$v_x$] (4,0);

        % 5. Fuente dependiente de tensión A*vx entre 'd' y 'a'
        % 'cV' para fuente de tensión controlada (rombo)
        \draw (d) to[cV, l=$A v_x$] (8,3) node[circ] {} node[above] {a};
        
        % 6. Resistencia R2 entre 'a' y 'b'
        \draw (8,3) to[R, l=$R_2$] (8,0) node[circ] {};
        
    \end{circuitikz}
\end{center}

\end{document}